\section{Cost Model and Break-Even}
\label{sec:cost}

Let a cooperative serve $A$ hectares with $N$ cropping seasons of device life,
and let the MRV cost per hectare per season decompose as
\begin{equation}
C=\underbrace{\frac{n_{\mathrm{dev}}\,p_{\mathrm{dev}}}{A\,N}}_{\text{hardware}}
+\underbrace{c_{\mathrm{conn}}}_{\text{connectivity}}
+\underbrace{h\,w}_{\text{labour}}
+\underbrace{\frac{c_{\mathrm{audit}}}{A}}_{\text{verification}},
\label{eq:cost}
\end{equation}
where $n_{\mathrm{dev}}$ and $p_{\mathrm{dev}}$ are the number and unit price of
devices, $h$ is the labour hours per hectare per season and $w$ the rural wage.
Revenue per hectare is
\begin{equation}
R=\kappa\,\rho ,
\label{eq:revenue}
\end{equation}
with $\kappa$ credits issued per hectare and $\rho$ the price per credit. The
scheme is economically viable when
\begin{equation}
C\ \le\ \beta R ,\qquad \beta\in(0,1),
\label{eq:breakeven}
\end{equation}
and $\beta=0.5$ is the level at which monitoring consumes half of the credit
revenue. Substituting \eqref{eq:revenue} gives the break-even issuance
\begin{equation}
\kappa_{\min}=\frac{C}{\beta\rho}.
\label{eq:kappamin}
\end{equation}

The tiers differ in $n_{\mathrm{dev}}$ and, more importantly, in which of the
three forgeable quantities of \eqref{eq:three} they can observe. Tier~0 is the
physics layer: $n_{\mathrm{dev}}=0$ beyond the controller, the log is written on
a handset the farmer already owns, and the filters of Sec.~\ref{sec:filters} run
on public weather data at no marginal cost. Tier~1 adds one flow meter at the
cooperative pumping station --- $n_{\mathrm{dev}}=1$ for the whole area $A$
rather than one per plot --- which observes the total volume. Tier~2 adds one
water-level node per cooperative, re-anchoring the twin and restoring the
calibration property a measurement-based MRV needs anyway.

None of the three observes timestamps, and by Prop.~\ref{prop:blind} the first
two are exactly invariant under the transformation that pays. The device anchor
of Sec.~\ref{sec:anchor} is therefore not an optional fourth tier but a
precondition of admissibility: it is produced by the same controller that
actuates the pump, so it adds storage and transmission of a record that already
exists, and no additional $n_{\mathrm{dev}}$.

Table~\ref{tab:baseline} compares the scheme with the three architectures it
replaces. With
$\kappa\in[4,6]$ and $\rho\in[\$15,\$25]$~\cite{rmi2025rice,upstream2026price},
revenue is $R\in[\$60,\$150]$ per hectare per season, so
$\kappa_{\min}$ at $\beta=0.5$ corresponds to a monitoring budget of
$\$30$--$75$ per hectare. The full stack sits an order of magnitude below that
ceiling, and the incremental cost of the anchor over a controller that is already
installed is about $1\%$ of conservative revenue. Table~\ref{tab:baseline}
decomposes the comparison against the three architectures it replaces; the
connectivity term is shown separately because it is the one that distinguishes a
sensor-per-plot design from a single shared gateway.

\begin{remark}
The device prices in Table~\ref{tab:baseline} are planning estimates, not
quotations. The ordering of the tiers and the break-even inequality
\eqref{eq:breakeven} do not depend on them; only the absolute margin does. A
quotation should replace the estimates before any procurement decision.
\end{remark}
